
\chapter{逻辑学}
\section{绪论}

逻辑学以思维形式及其规律为研究对象。它并不直接考察具体学科对象本身，而是考察各种思想内容在抽去个别题材之后所呈现出来的结构，以及这些结构在不同代入下表现出来的真假特征。正因为如此，逻辑学能够为判定推理是否在形式上正确提供一般标准。

\subsection{逻辑学的对象}

在日常思维、科学论证与数学推演中，人们所处理的题材可以彼此不同，但其中往往包含相同的思维结构。逻辑学所抽取和研究的，正是这种不依赖具体题材而可以重复出现的结构性方面，以及由这种结构所体现出来的规律性。

\paragraph{思维形式}\label{logic.thought_form}

思维形式是思维内容的存在方式和联系方式。它是在撇开具体对象与个别事实之后，对思想结构所做的逻辑抽象。因此，不同的思想内容可以具有同一思维形式，而同一思维形式也可以在不同领域中反复出现。

例如，下列语句虽然分别讨论植物、动物和经济对象，但它们具有共同的结构：

\[
	\text{所有 } S \text{ 是 } P.
\]

同样地，一类条件陈述也可以抽象为如下形式：

\[
	\text{如果 } p\text{，那么 } q.
\]

前一形式刻画词项之间的归属关系，后一形式刻画命题之间的条件联系。逻辑学正是通过这种抽象，把纷繁的具体内容转化为可分析的形式结构。

\paragraph{逻辑常项与逻辑变项}\label{logic.logical_constant}\label{logic.logical_variable}

思维形式由逻辑常项和逻辑变项构成。逻辑常项是形式中的不变成分，用来规定思想的结构和连接方式；逻辑变项是形式中的可变成分，用来容纳具体内容。就上面的两个例子而言，“所有\dots 是\dots”和“如果\dots 那么\dots”属于逻辑常项，而 $S$、$P$、$p$、$q$ 则属于逻辑变项。

词项变项可以由不同的概念或词项代入，命题变项则可以由不同的命题代入。于是，相同的逻辑结构能够承载不同学科与不同语境中的具体思想内容，而逻辑分析所关注的正是这种结构在代入变化之下仍保持不变的形式方面。

\begin{note}
逻辑常项决定了一个表达式属于何种形式，例如全称判断形式或条件命题形式；逻辑变项则使这一形式能够应用到具体对象和具体断言之中。二者结合起来，才构成完整的思维形式。
\end{note}

\paragraph{逻辑规律与逻辑矛盾}\label{logic.logical_law}\label{logic.logical_contradiction}

思维形式本身不带有特定的经验内容，因此单独看时并不直接表达某个具体真或假的判断。只有当具体词项或命题代入逻辑变项之后，这一形式才被赋予确定内容，并进而呈现真假。逻辑学特别关心的是那些在任意代入之下都保持同一真假特征的形式。

如果一个思维形式无论怎样代入都表达真实内容，那么它就是逻辑规律。例如：

\[
	\text{所有 } S \text{ 是 } S, \qquad p \lor \neg p.
\]

如果一个思维形式无论怎样代入都表达虚假内容，那么它就是逻辑矛盾。例如：

\[
	\text{有 } S \text{ 不是 } S, \qquad p \land \neg p.
\]

逻辑规律给出形式上必真的结构，逻辑矛盾给出形式上必假的结构。逻辑学一方面利用逻辑规律作为有效推理的依据，另一方面排除逻辑矛盾，以保证思维在形式上的一致性与正确性。

\begin{important}
研究对象可以概括为三个彼此关联的层面：其一，思维形式本身；其二，逻辑常项与逻辑变项所构成的结构；其三，这些结构在任意代入下所表现出的逻辑规律与逻辑矛盾。逻辑学正是在这一基础上判定推理的形式有效性。
\end{important}

\subsection{思维、语言和逻辑}

思维、语言和逻辑彼此区分，但不能彼此割裂。思维是人类理性认识的基本活动，语言是思维获得直接现实性的物质外壳，逻辑学则通过分析语言的表达形式来把握思维的形式结构及其规律。因此，逻辑学既不等同于语言学，也不脱离语言而直接把握不可见的思维，而是借助语言这一中介来研究正确思维的形式条件。

\paragraph{思维、语言和逻辑的关系}\label{logic.relation_between_thought_language_and_logic}

思维以概念、判断和推理等基本形式反映对象世界，但思维本身看不见、听不到，只有通过语言表达出来，才获得可以被分析和交流的现实形态。语言因而不是外加于思维的偶然工具，而是思维进入公共讨论、论证和检验过程的必要媒介。

对于逻辑学而言，这一点尤其重要。逻辑学所研究的是思维形式，而思维形式总是通过某种语言表达出来。于是，逻辑学通常并不直接分析个体心理活动，而是透过语词、语句和论证结构来考察其中所体现的概念关系、命题结构和推理形式。

\begin{note}
可以把三者的关系概括为：思维提供逻辑研究的对象内容，语言提供逻辑研究的现实载体，逻辑则从语言表达中抽取和评价思维的形式结构。
\end{note}

\paragraph{对象语言与元语言}\label{logic.object_language_and_metalanguage}

既然逻辑学通过语言研究思维，就必须区分被研究的语言和用于研究的语言。前者称为对象语言，后者称为元语言。对象语言是逻辑学直接分析的表达系统；元语言则用来描述对象语言的符号、规则、结构和语义。

在具体场景中，对象语言和元语言可以是不同语言。例如，中国人用汉语解释英语语法时，英语是对象语言，汉语是元语言。它们也可以是同一种语言在不同层次上的使用；例如，用汉语来讨论汉语命题的结构时，被分析的汉语表达式是对象语言，而说明这些表达式如何构成和如何解释的汉语，则充当元语言。

区分对象语言和元语言，有助于避免把语言中的表达和关于该表达的说明混为一谈，这也是后续形式语义与元理论分析的基础。

\paragraph{自然语言}\label{logic.natural_language}

自然语言是人类在长期社会实践中形成并约定俗成地使用的语言系统，例如汉语、英语和日语。它首先服务于日常思维、社会交往和经验表达，因此具有丰富、灵活和适应语境的优点。

从逻辑分析的角度看，语言至少包含三个基本方面：

\begin{enumerate}
	\item 基本符号，即构成语言的基本材料；
	\item 语形规则，即规定哪些符号串是合式表达；
	\item 语义规则，即规定合式表达具有什么意义。
\end{enumerate}

自然语言的突出特点在于它通常带有歧义，并且其确切意义常常依赖语境、说话者意图和具体交际环境。正因为自然语言既强大又含混，逻辑学一方面离不开它作为对象语言和日常推理的基本载体，另一方面又需要借助更严格的表达工具来减少歧义。

\paragraph{符号语言与形式系统}\label{logic.symbolic_language}\label{logic.formal_system}

为适应严格研究的需要，人们构造出人工语言。逻辑学中使用的人工语言通常称为符号语言。与自然语言相比，符号语言通过明确规定基本符号、构造规则与解释规则，尽可能排除歧义，使表达式的结构和推理步骤得到精确控制。

当这种符号语言被进一步严格化时，就形成形式语言；以形式语言为基础，并把合式公式、推导规则以及必要的解释方式明确规定下来，就得到形式系统。形式系统的意义在于，它把推理从日常语感中分离出来，使人们能够在抽象层面精确研究结论是否由前提必然推出。

因此，形式系统不是脱离思维现实的纯技术装置，而是对思维形式进行严格刻画的工具。现代逻辑中的演算系统、证明系统和语义解释框架，都建立在这一思想之上。

\paragraph{传统逻辑与现代逻辑}\label{logic.traditional_and_modern_logic}

传统逻辑与现代逻辑在总目标上具有连续性，二者都以研究正确思维、分析论证和判定推理有效性为中心。二者的主要差别，不在于是否研究思维，而在于所使用的工具语言及其形式化程度。

传统逻辑主要借助自然语言来分析概念、判断和推理，因而在表达上更贴近日常论证。现代逻辑则更多依赖符号语言和形式化方法，把命题结构、量化结构和证明规则精确地表示出来，从而能够更系统地研究逻辑系统的性质。

在这一发展过程中，一阶逻辑构成现代逻辑的基础部分。它并不是对传统逻辑研究对象的否定，而是对传统逻辑中命题和词项推理研究的精确化、严格化与扩展。也正因此，理解自然语言与符号语言、对象语言与元语言之间的关系，是进入现代逻辑的必要准备。

\begin{important}
这一子节的核心线索是：思维需要通过语言取得可分析的现实形式，逻辑学需要通过分析语言来把握思维形式；自然语言提供原始的思想材料，符号语言和形式系统提供严格的分析工具，而现代逻辑正是在这种工具自觉的基础上发展起来的。
\end{important}

\subsection{逻辑学的性质和作用}

在前两节中，逻辑学已经被界定为研究思维形式及其规律的学科。进一步说，逻辑学不仅有自己的研究对象，而且具有明确的学科性质，并在认识、论证和交际活动中发挥持续作用。教材把这种性质概括为基础性、工具性和全人类性，并进一步强调逻辑学在探求真理、辨析论证和提高表达质量方面的普遍意义。

\paragraph{逻辑学的基础性}\label{logic.foundational_character_of_logic}

逻辑学具有基础性，因为只要一个学科涉及概念、判断和推理，就离不开逻辑。无论是数学证明、自然科学解释，还是社会科学分析与日常议论，都会使用概念来把握对象、使用判断来陈述结论、使用推理来建立前提与结论之间的联系。

正由于逻辑学研究这些最一般的思维形式及其规律，它并不局限于某一门专门学科，而是构成多种知识活动共同依赖的方法基础。从这个意义上说，逻辑学并不是与其他学科并列的一种狭窄专业知识，而是许多学科训练中都需要反复运用的基础能力。

\paragraph{逻辑学的工具性}\label{logic.instrumental_character_of_logic}

逻辑学同时具有工具性。它为人们提供的不只是关于思维形式的理论说明，更是一套可实际运用的分析工具和评价标准。借助逻辑，人们能够检查一个论证是否有效，辨认一个判断是否含混，分析一个概念是否界定清楚，也能够更有条理地组织自己的表达。

这种工具性体现在两个方向上。一方面，逻辑帮助人们形成更规范的思维过程，使结论的推出更有根据；另一方面，逻辑也帮助人们改进交际与论辩，使表达更准确、回应更切题、驳斥更有力。因此，逻辑既服务于认识活动，也服务于论证和交流活动。

\paragraph{逻辑学的全人类性}\label{logic.universality_of_logic}

逻辑学还具有全人类性。不同民族、不同地区和不同文化传统在语言、习俗和思想资源上可以存在巨大差异，但只要讨论的是有效推理、形式一致性和正确论证，就会诉诸某种普遍适用的共同逻辑。逻辑并不以某一种民族语言或某一种地方文化为专属前提。

这并不否认不同文明中逻辑研究的历史形式各不相同，而是强调：尽管表达方式和发展路径可以不同，但正确推理所依赖的形式要求具有普遍性。也正因此，逻辑学能够成为跨文化的共同知识，并构成全人类认识和论证活动的共有精神资源。

\begin{note}
基础性说明逻辑学在知识体系中的地位，工具性说明逻辑学在实际思维与表达中的使用方式，全人类性则说明逻辑学的适用范围并不受民族语言和地域文化的限制。
\end{note}

\paragraph{逻辑学的作用}\label{logic.function_of_logic}

逻辑学的作用，集中体现为它能够以逻辑形式与逻辑规律的知识来约束和规范认识活动与交际活动。人们在获取新知、进行论证、形成判断和反驳谬误时，并不是随意联想，而是需要依靠前提与结论之间合理的形式联系。逻辑学正是在这一点上为思维活动提供方法自觉。

具体地说，逻辑学至少有三方面作用。第一，它有助于提高表达的准确性，使概念使用、命题陈述和论证结构更清晰。第二，它有助于提高论证的有效性，使人们能够区分有效推理、无效推理和似是而非的诡辩。第三，它有助于培养辨识谬误和探求真理的能力，使思维不至于陷入混乱、矛盾或无根据的武断。

因此，学习逻辑学的意义并不止于掌握若干术语或规则，而在于形成一种对概念、判断和推理负责的思维习惯。它既是学术研究的基础训练，也是公共表达、理性讨论和日常决策的重要支持。

\begin{important}
这一子节的要点在于：逻辑学之所以重要，不仅因为它研究思维形式，更因为它在知识体系中具有基础地位，在实际思维与论证中具有工具价值，并且作为关于正确推理的学问，对全人类都具有普遍适用性。
\end{important}
